\chapter{Il questionario: la webapp di raccolta dati}
\label{ch:questionario}

Tutta l'informazione a priori che rende trattabile la disaggregazione arriva dal questionario.
Per raccoglierla è stata sviluppata una \textbf{applicazione web dedicata}, che sostituisce la
compilazione manuale su foglio di calcolo usata in precedenza. Questo capitolo descrive che cosa
viene chiesto al cliente, come è organizzata l'intervista e in che forma i dati arrivano al
modello di ottimizzazione.

\section{Perché una webapp e non un modulo statico}

Il questionario precedente era un elenco fisso di elettrodomestici: una riga per
\emph{tipologia}, con il numero di esemplari posseduti. Tre limiti lo rendevano insufficiente.

\begin{center}
\begin{tabularx}{\textwidth}{@{}l X@{}}
\toprule
\textbf{Limite} & \textbf{Conseguenza sul modello} \\
\midrule
Elenco chiuso di tipologie
& Un apparecchio non previsto dall'elenco non era dichiarabile e finiva nel rumore $\eta(t)$
  della~\eqref{eq:modello-misura}. \\
\addlinespace
Una riga per tipologia
& Due climatizzatori erano costretti a condividere fasce orarie e consumi, anche se uno sta in
  camera e l'altro in salotto e vengono usati in momenti opposti della giornata. \\
\addlinespace
Nessuna potenza dichiarata
& La potenza $P_i$ veniva presa da valori di letteratura per tipologia, senza alcuna
  possibilità di correzione sul singolo apparecchio. \\
\addlinespace
Domande non utilizzate
& Superficie, anno di costruzione, tipo di riscaldamento, numero di occupanti: dati raccolti
  ma mai letti da nessun modello, che allungavano l'intervista senza aggiungere informazione. \\
\bottomrule
\end{tabularx}
\end{center}

La webapp risolve i primi tre punti e rimuove il quarto: l'intervista è più corta e chiede
solo grandezze che entrano effettivamente nella funzione obiettivo.

\begin{notebox}{Il cambio strutturale: istanze, non tipologie}
Nel nuovo schema l'elenco dei dispositivi è una \textbf{lista di istanze indipendenti}: ogni
apparecchio fisico presente in casa è una voce a sé, con la propria etichetta
(\code{Camera}, \code{Garage}), le proprie potenze e i propri orari. Un'abitazione con due
frigoriferi e tre climatizzatori dichiara cinque istanze distinte, non due righe con un contatore.
\end{notebox}

\section{Come si svolge l'intervista}

L'applicazione è realizzata in Python con Streamlit (\code{questionnaire/app.py}) e si compila da
browser, senza installazione lato cliente. Il flusso è in tre fasi.

\subsection*{Fase 1 --- Dati della casa}

Si identifica l'abitazione tramite l'\code{IMEI} del dispositivo IoT che ne misura il consumo e si
raccolgono due sole grandezze:

\begin{itemize}[nosep,leftmargin=1.4em]
  \item la \textbf{potenza contrattuale} in kW, che fissa un tetto fisico alla somma degli
        assorbimenti simultanei;
  \item il \textbf{carico base} in W, cioè quanto assorbe la casa quando tutto sembra spento
        (router, allarme, standby, illuminazione di servizio), con il relativo livello di
        confidenza.
\end{itemize}


\subsection*{Fase 2 --- Aggiunta dei dispositivi}

Il cliente aggiunge un dispositivo alla volta, scegliendo la tipologia da un \textbf{catalogo}
di 18 tipologie, organizzate per categoria: freddo, lavaggio, cottura, climatizzazione,
intrattenimento, mobilità. Per ogni istanza si raccolgono:

\begin{center}
\begin{tabularx}{\textwidth}{@{}l X@{}}
\toprule
\textbf{Campo} & \textbf{Significato} \\
\midrule
Etichetta & Nome libero che distingue due apparecchi della stessa tipologia. \\
\addlinespace
Marca e modello (opz.)
& Selezionandoli dal catalogo, le potenze vengono \textbf{precompilate} automaticamente; restano
  comunque modificabili a mano. \\
\addlinespace
Potenza tipica, minima, di picco
& I tre valori in W dell'assorbimento in funzione. La potenza tipica è il $P_i$ del modello, i due
  estremi definiscono l'intervallo entro cui il solver può muovere il livello di assorbimento. \\
\addlinespace
Modalità d'uso
& \code{cycle}, \code{daily\_hours} o \code{always\_on}: decide quali domande di utilizzo vengono
  mostrate (tabella seguente). \\
\addlinespace
Fascia oraria
& Ora di inizio e di fine della finestra in cui l'apparecchio viene tipicamente avviato o usato. \\
\addlinespace
Mesi di attività
& I mesi dell'anno in cui l'apparecchio viene usato: è la dichiarazione di
  \textbf{stagionalità} (un climatizzatore attivo solo da giugno a settembre, una stufa solo
  d'inverno). \\
\addlinespace
Confidenza
& \code{alta}, \code{media} o \code{bassa}: quanto il cliente è sicuro di ciò che ha dichiarato
  per quell'apparecchio. \\
\bottomrule
\end{tabularx}
\end{center}

Le domande sull'utilizzo non sono le stesse per tutti: chiedere a un frigorifero
\emph{«quante volte alla settimana lo usi?»} non ha senso. Ogni tipologia del catalogo porta con sé
una \textbf{modalità d'uso} che seleziona il gruppo di domande pertinenti.

\begin{center}
\begin{tabularx}{\textwidth}{@{}l X l@{}}
\toprule
\textbf{Modalità} & \textbf{Domande poste} & \textbf{Esempi} \\
\midrule
\code{cycle}
& Frequenza settimanale (min--max), durata del ciclo in minuti (min--max), fascia oraria di avvio
& lavatrice, forno, auto elettrica \\
\addlinespace
\code{daily\_hours}
& Ore medie di accensione al giorno (min--max), fascia oraria di utilizzo
& climatizzatore, TV, computer \\
\addlinespace
\code{always\_on}
& Nessuna domanda di utilizzo: solo potenza media e duty cycle
& frigorifero, congelatore \\
\bottomrule
\end{tabularx}
\end{center}

Tutte le grandezze di utilizzo sono chieste come \textbf{intervallo minimo--massimo}, non come
valore puntuale. È coerente con il principio guida del capitolo precedente: il cliente risponde a
memoria, e un intervallo dichiara onestamente la propria incertezza dove un numero secco fingerebbe
una precisione che non esiste.

\begin{notebox}{Fasce a cavallo della mezzanotte}
Se l'ora di fine precede quella di inizio --- un climatizzatore usato dalle 23:00 alle 07:00 ---
l'applicazione riconosce il caso e marca la finestra come estesa al giorno successivo, invece di
trattarla come un intervallo vuoto.
\end{notebox}

\subsection*{Fase 3 --- Validazione ed export}

Prima del salvataggio la risposta viene controllata (\code{questionnaire/schema.py}): IMEI
presente e numerico, potenza contrattuale positiva, almeno un dispositivo, coerenza degli
intervalli (la potenza minima non può superare la tipica, il minimo di ogni intervallo non può
superare il massimo, le ore giornaliere non possono superare 24), etichette non duplicate
all'interno della stessa tipologia. I problemi vengono mostrati in chiaro al compilatore, che
corregge prima di inviare.

L'export produce due formati dello stesso contenuto: un \textbf{JSON} in schema v3, che è il
formato di riferimento, e un \textbf{Excel} a due fogli (\code{household} con i dati casa,
\code{devices} con una riga per istanza) per l'ispezione manuale.

\section{Il formato dei dati prodotti}

Un estratto della struttura JSON, con due climatizzatori dichiarati separatamente:

{\small
\begin{verbatim}
{
  "schema_version": 3,
  "imei": "860123456789012",
  "household": {
    "contract_power_kw": 4.5,
    "base_load_w": 85.0,
    "base_load_confidence": "media"
  },
  "devices": [
    {"id": "climatizzatore-a1b2c3", "type": "Climatizzatore", "label": "Camera",
     "usage_mode": "daily_hours",
     "daily_usage_hours_min": 6, "daily_usage_hours_max": 8,
     "start_window_start_min": 1380, "start_window_end_min": 420,
     "start_window_wraps_next_day": true,
     "active_months": [6, 7, 8, 9],
     "p_min_w": 400.0, "p_typical_w": 1100.0, "p_max_w": 1600.0,
     "confidence": "media"},
    {"id": "climatizzatore-d4e5f6", "type": "Climatizzatore", "label": "Salotto",
     "usage_mode": "daily_hours",
     "daily_usage_hours_min": 2, "daily_usage_hours_max": 4, "...": "..."}
  ]
}
\end{verbatim}
}

\noindent
Gli orari sono espressi in minuti dalla mezzanotte (1380 $=$ 23:00, 420 $=$ 07:00), i mesi come
numeri da 1 a 12.

\section{Dove finisce ciascuna risposta}

Ogni campo raccolto ha una destinazione precisa nel modello. La tabella seguente è la mappa fra
intervista e formulazione matematica.

\begin{center}
\begin{tabularx}{\textwidth}{@{}l l X@{}}
\toprule
\textbf{Risposta} & \textbf{Simbolo} & \textbf{Uso nel modello} \\
\midrule
Presenza del dispositivo & --- & Determina l'insieme degli apparecchi modellati: un
  apparecchio non dichiarato non ha variabili. \\
\addlinespace
Potenza tipica & $P_i$ & Ampiezza del gradino associato all'apparecchio
  (capitolo~\ref{ch:modelli}). \\
\addlinespace
Potenza min./picco & $\delta$ & Delimitano i livelli di assorbimento alternativi che il solver può
  scegliere. \\
\addlinespace
Ore giornaliere / durata ciclo & $E_i$ & Monte ore giornaliero atteso; lo scostamento è penalizzato
  con $\lambda_d$ (capitolo~\ref{ch:obiettivo}). \\
\addlinespace
Frequenza settimanale & --- & Quota settimanale di accensioni, propagata da un giorno all'altro
  come costante (capitolo~\ref{ch:pipeline}). \\
\addlinespace
Fascia oraria & $f_i(t)$ & Fattore temporale: penalità crescente con la distanza dalla finestra
  dichiarata (capitolo~\ref{ch:fattore}). \\
\addlinespace
Mesi di attività & $\lambda_s$ & Penalità di stagionalità per ogni slot acceso fuori dai mesi
  dichiarati. \\
\addlinespace
Modalità \code{always\_on} & $p^{\text{ao}}_j$ & L'apparecchio entra come livello di assorbimento
  continuo ottimizzato dal solver, senza attivazioni discrete. \\
\addlinespace
Carico base della casa & --- & Quota costante non attribuibile ad alcun apparecchio, che evita di
  gonfiare le stime dei sempre accesi. \\
\addlinespace
Potenza contrattuale & --- & Tetto fisico alla somma degli assorbimenti simultanei. \\
\addlinespace
Confidenza & --- & Destinata a modulare per apparecchio il peso delle penalità: una dichiarazione
  a confidenza bassa deve costare meno da violare. \\
\bottomrule
\end{tabularx}
\end{center}

\begin{warnbox}{Stato di integrazione}
La webapp è operativa e produce dati in schema v3. Restano due passi per l'uso in produzione:
\begin{itemize}[nosep,leftmargin=1.4em]
  \item il \textbf{catalogo modelli} è al momento uno stub dimostrativo, con marca generica e
        potenze indicative; va sostituito con un database prodotti reale (etichette energetiche
        europee, API dei produttori). L'applicazione dipende solo dalla firma della funzione di
        ricerca, quindi la sostituzione è circoscritta;
  \item serve l'\textbf{adattatore} che converta le istanze v3 nei profili letti dal modello di
        ottimizzazione, che oggi consuma ancora il formato v2.
\end{itemize}
Il peso della confidenza sulle penalità è previsto dallo schema ma non ancora letto dal modello.
\end{warnbox}
